\mxnewpage\section*{Correction EV1-T1}\label{EV1-T1-CORR}
\textbf{QCM :} 1 B, 2 C, 3 A, 4 C, 5 D, 6 D, 7 B, 8 D, 9 D, 10 B, 11 B, 12 A. Chaque réponse correcte : \num{0,5} point ; total 6.
\paragraph{Vérifications des automatismes.} Priorités ; somme de fractions ; soustraction des exposants ; coefficient multiplicateur \num{1,2} ; division par \num{0,8} ; produit \num{1,25} fois \num{0,8} ; résolution affine ; accroissement vertical divisé par l’horizontal ; moyenne ; angle divisé par 360 ; complément à 1 ; substitution entre parenthèses.
\subsection*{Exercice 1 -- Un fonds de livres}
\begin{enumerate}
\item $u_n=120+18n$, arithmétique de raison 18. \textit{Barème : 1 point. Résultat et justification élémentaire attendus ; demi-point possible pour démarche correcte inachevée.}
\item $u_8=264$ ; neuf termes, somme $\dfrac{9(120+264)}{2}=1728$. Cette somme cumule des états mensuels, pas des livres distincts. \textit{Barème : 2 points. Méthode ou justification : 1 ; calcul et conclusion : 1.}
\item $120+18n>300\iff n>10$. Premier rang 11 ; $u_{10}=300$, $u_{11}=318$. \textit{Barème : 2 points. Méthode ou justification : 1 ; calcul et conclusion : 1.}
\end{enumerate}
\subsection*{Exercice 2 -- Aire et choix de forme}
\begin{enumerate}
\item $A=-x^2+14x=49-(x-7)^2$. \textit{Barème : 1 point. Résultat et justification élémentaire attendus ; demi-point possible pour démarche correcte inachevée.}
\item $A\leq49$, égalité en 7. \textit{Barème : 2 points. Méthode ou justification : 1 ; calcul et conclusion : 1.}
\item $49-(x-7)^2\geq40\iff(x-7)^2\leq9$, donc $\intff{4}{10}$. \textit{Barème : 2 points. Méthode ou justification : 1 ; calcul et conclusion : 1.}
\end{enumerate}
\subsection*{Exercice 3 -- Un angle droit}
\begin{enumerate}
\item $\vect{AB}=\pt{3;1}$, $\vect{AC}=\pt{-1;3}$. \textit{Barème : 1 point. Résultat et justification élémentaire attendus ; demi-point possible pour démarche correcte inachevée.}
\item Produit $-3+3=0$. \textit{Barème : 1 point. Résultat et justification élémentaire attendus ; demi-point possible pour démarche correcte inachevée.}
\item $BC=\sqrt{(-4)^2+2^2}=2\sqrt5$ ; côtés perpendiculaires de norme $\sqrt{10}$, aire 5. \textit{Barème : 2 points. Méthode ou justification : 1 ; calcul et conclusion : 1.}
\end{enumerate}
\paragraph{Étalonnage.} Durée cible : 20–25 min pour les automatismes, 80–85 min pour les exercices, 10–15 min de relecture. Ces durées sont des estimations éditoriales à mesurer en classe. Accepter toute méthode correcte ; ne pas sanctionner deux fois une erreur numérique propagée avec cohérence. Le barème proposé ne vaut pas décision d’harmonisation de l’équipe.
